\section{Conclusion and Future Work}
\label{sec:conclusion}

This paper presented a reliability-aware, policy-constrained self-healing control plane for data pipelines. The framework models self-healing as a staged reliability lifecycle rather than as a single detection or retry mechanism. The lifecycle includes failure detection, root-cause classification, remediation selection, policy approval, execution, post-recovery validation, rollback or escalation, and evidence preservation.

The evaluation used a reproducible synthetic failure-injection setup with 780 total trials, including injected failures, healthy controls, and boundary controls. Under the evaluated controlled experiment configuration, the framework achieved 100\% failure detection accuracy and 100\% root-cause classification accuracy. Verified recovery was achieved in 52.17\% of injected failure trials. Runtime was measured across 780 records, with a mean self-healing cycle runtime of 12.609 ms, a median of 10.609 ms, and p95 of 26.666 ms.

These results support the feasibility of the proposed framework under controlled experimental conditions. More importantly, they show why detection, diagnosis, and verified recovery should be evaluated separately. Accurate detection does not guarantee safe remediation. Correct classification does not guarantee that a pipeline should continue. A recovery action should be considered successful only when it is policy-approved, executed, and validated after execution.

The main conclusion is that self-healing data pipelines should be treated as governed reliability-control systems. A useful framework must not only identify failures, but also determine whether automated remediation is safe, whether invalid output propagation should be blocked, and whether sufficient evidence has been preserved for audit and later inspection. This distinction is especially important in enterprise and regulated data environments, where unrestricted automation can create operational or compliance risk.

The study also has clear boundaries. The evaluation uses synthetic data, controlled failure injection, and a local research execution environment. These choices improve reproducibility but limit direct generalization to all production pipelines. The results should therefore be interpreted as evidence of architectural and experimental feasibility rather than proof of universal production readiness.

Future work should extend the framework in several directions. First, the evaluation should be repeated using production-like datasets, historical incident traces, and larger workloads. Second, the failure taxonomy should be expanded to include compound failures, intermittent upstream outages, delayed data, permission changes, and downstream contract violations. Third, the framework should be compared against alert-only, retry-only, and static-rule baselines under identical scenarios. Fourth, future versions should explore probabilistic, confidence-aware, hierarchical, or multi-label root-cause diagnosis. Finally, production-oriented work should study integration with orchestration platforms, observability systems, policy-as-code engines, incident-management workflows, and enterprise governance controls.

Overall, the paper contributes a reproducible implementation and evaluation of an end-to-end self-healing control loop for data pipelines. The work demonstrates that self-healing claims are strongest when they are supported by separate evidence for detection, classification, recovery, runtime behavior, and auditability.
